% !TeX document-id = {1c0b4298-276a-4038-8e08-c4a1f4846da7}
% !TeX TXS-program:bibliography = txs:///biber
\documentclass[12pt,a4paper]{article}
\usepackage[T1]{fontenc}
\usepackage{graphicx}
\usepackage{mathtools}
\usepackage{amssymb}
\usepackage{amsthm}
\usepackage{thmtools}
\usepackage{xcolor}
\usepackage{nameref}
\usepackage{hyperref}
\usepackage{color}
\usepackage{float}
\newtheorem{theorem}{Teorema}
%\usepackage[margin=2.5cm]{geometry}
\usepackage[brazilian]{babel}
\usepackage[
backend=biber,
style=authoryear,
natbib=true
]{biblatex}
\addbibresource{../bibliography.bib}
\newtheorem{definition}{Definição}
\newtheorem{lemma}{Lema}
\newtheorem{proposition}{Proposição}
\title{Inferência no modelo linear Gaussiano}
\author{ Professor Luís Antonio Fantozzi Alvarez}
\date{15 de abril de 2026}
\begin{document}
	\maketitle
	O objetivo destas notas é verificar que, no modelo linear Gaussiano:
	
	$$\boldsymbol{y} = \boldsymbol{X}\beta + \boldsymbol{\epsilon}\, ,  \beta \in \mathbb{R}\, ,$$
	$$\operatorname{posto}(\boldsymbol{X}) = k\, ,$$
	$$\boldsymbol{\epsilon}|\boldsymbol{X} \sim \mathcal{N}(0_{n\times 1},\sigma^2 \mathbb{I}_{n\times n})\, ,$$
	dada uma matriz $r \times k$ $R$ de posto $r$ e um vetor $r\times 1$ $c$ , a estatística de teste:
	
	$$\footnotesize W_{R,c} = \frac{(R\hat{b}-c)'\left(R\hat{\sigma}^2(\boldsymbol{X}'\boldsymbol{X})^{-1}R'\right)^{-1}(R\hat{b}-c)}{r}  = \frac{(R\hat{b}-c)'\left(R(\boldsymbol{X}'\boldsymbol{X})^{-1}R'\right)^{-1}(R\hat{b}-c)}{r \hat{\sigma}^2}\, ,$$
	onde $\hat{\sigma}^2 = \frac{1}{n-k}\sum_{i=1}^n (Y_i -X_i'\hat{b})^2$, terá distribuição, sob a hipótese nula $R\beta = c$,
	
	$$W_{R,c}|\boldsymbol{X}\sim F(r,n-k)$$
	
	Dividimos a demonstração em quatro etapas:
	
\paragraph{Etapa 0} Para uma matriz $A$ simétrica, a decomposição espectral permite-nos escrever $A = P \Lambda P'$, onde $\Lambda$ é uma matriz diagonal com os autovalores de $A$, e $P$ é uma matriz ortogonal ($P'P=PP'=\mathbb{I}$). Se $A$ é \textbf{positiva definida}, sabemos que os autovalores de $A$ são estritiamente positivos e, nesse caso, definimos a potência $-1/2$ de $A$ como:

$$A^{-1/2}:= P\Lambda^{-1/2}P'\, .$$

Observe que, com base nessa definição, $$A^{-1/2}A^{-1/2} =  P\Lambda^{-1/2}\underbrace{P' P}_{=\mathbb{I}}\Lambda^{-1/2}P' = P\Lambda^{-1}P'=A^{-1}\, .$$

\paragraph{Etapa 1} Defina $N := (R\hat{b}-c)'\left(R\sigma^2(\boldsymbol{X}'\boldsymbol{X})^{-1}R'\right)^{-1}(R\hat{b}-c) $. Observe que, sob a hipótese nula, $R\hat{b}-c|\boldsymbol{X}\sim \mathcal{N}(0_{r \times 1}, \sigma^2 R(\mathbf{X}'\mathbf{X})^{-1}R')$. Como $R(\mathbf{X}'\mathbf{X})^{-1}R'$ é positiva definida pela condição de posto sobre $R$,\footnote{Como $R$ tem posto $r$, para todo $s \in \mathbb{R}^r$, $R's\neq 0$. Consequentemente, como $(\mathbf{X}'\mathbf{X})^{-1}$ é positiva definida (use a decomposição espectral e o fato de que $\mathbf{X}'\mathbf{X}$ é positiva definida para verificar isso), temos que: $s' R(\mathbf{X}'\mathbf{X})^{-1}R's=(R's)'(\mathbf{X}'\mathbf{X})^{-1}(R's) > 0$.}, sua inversa de fato existe, e temos que, sob a nula:
$$(\sigma^2 R(\mathbf{X}'\mathbf{X})^{-1}R')^{-1/2}(R\hat{b}-c) |\boldsymbol{X}\sim \mathcal{N}(0_{r\times1},\mathbb{I}_{r\times r})\, .$$
Consequentemente, condicionalmente a $\mathbf{X}$, $N = \lVert (\sigma^2 R(\mathbf{X}'\mathbf{X})^{-1}R')^{-1/2}(R\hat{b}-c) \rVert_2^2$ é a soma dos quadrados de $r$ normais padrão independentes, de modo que $N|\mathbf{X}\sim \chi^2(r)$.

\paragraph{Etapa 2} Defina $D:=(n-k)\hat{\sigma}^2/\sigma^2$. Observe que, $\sigma^2D = \lVert M\boldsymbol{y} \rVert_2^2 \overset{M\boldsymbol{X}=0_{n\times k}}{=} \lVert M\boldsymbol{\epsilon} \rVert_2^2= (M\boldsymbol{\epsilon})'(M\boldsymbol{\epsilon}) = \boldsymbol{\epsilon}'M\boldsymbol{\epsilon}$. Pela decomposição espectral, podemos escrever $M = P \Lambda P'$, onde $PP'=P'P=\mathbb{I}_{n\times n}$. Dessa forma, obtemos que:

$$D = \left(\sigma^{-1}P\boldsymbol{\epsilon}\right)'\Lambda \left(\sigma^{-1}P\boldsymbol{\epsilon}\right)\, ,$$
onde $\sigma^{-1}P\boldsymbol{\epsilon}|\boldsymbol{X}\sim \mathcal{N}(0_{n\times n}, \mathbb{I}_{n\times n})$. Consequentemente, como $\Lambda$ é uma matriz cuja diagonal principal consiste de zeros ou uns, com $1$ aparecendo $n-k$ vezes, temos que $D$, condicionalmente a $\boldsymbol{X}$, se escreve como a soma dos quadrados de $n-k$ normais padrão independentes. Consequentemente, $D|\boldsymbol{X}\sim \chi^2(n-k)$.

\paragraph{Etapa 3} Vamos mostrar, agora, que $N$ e $D$ são independentes, condicionalmente a $\boldsymbol{X}$. Observe que, sob a nula e condicionalmente a $\boldsymbol{X}$, $D$ e $N$ são transformações, respectivamente, das variáveis aleatórias $R\hat{b}-c=R(\boldsymbol{X}'\boldsymbol{X})^{-1}\boldsymbol{X}'\boldsymbol{\epsilon}$ e $M\boldsymbol{\epsilon}$. Consequentemente, é suficiente mostrar que $R(\boldsymbol{X}'\boldsymbol{X})^{-1}\boldsymbol{X}'\boldsymbol{\epsilon}$ e $M\boldsymbol{\epsilon}$ são independentes, condicionalmente a $\boldsymbol{X}$. Como $\begin{bmatrix}
R(\boldsymbol{X}'\boldsymbol{X})^{-1}\boldsymbol{X}' \\
M
\end{bmatrix}\boldsymbol{\epsilon}$ tem distribuição normal multivariada (condicional a $\boldsymbol{X}$), é suficiente verificar que $\operatorname{cov}(R(\boldsymbol{X}'\boldsymbol{X})^{-1}\boldsymbol{X}'\boldsymbol{\epsilon},M\boldsymbol{\epsilon}|\boldsymbol{X}) =0_{r \times n}$. De fato:

\begin{eqnarray*}
	\begin{aligned}
		\operatorname{cov}(R(\boldsymbol{X}'\boldsymbol{X})^{-1}\boldsymbol{X}'\boldsymbol{\epsilon},M\boldsymbol{\epsilon}|\boldsymbol{X}) =\\ \mathbb{E}[R(\boldsymbol{X}'\boldsymbol{X})^{-1}\boldsymbol{X}'\boldsymbol{\epsilon}\boldsymbol{\epsilon}'M|\boldsymbol{X}]=\sigma^2 R\boldsymbol{X}'(\boldsymbol{X}'\boldsymbol{X})^{-1}\underbrace{\boldsymbol{X}'M}={0_{k \times n}} = 0_{r \times n}
	\end{aligned}
\end{eqnarray*}
de onde concluímos o resultado desejado.

\paragraph{Etapa 4} Para concluir, observamos que:
$$\footnotesize W_{R,c} =\frac{N}{D}\frac{n-k}{r}\, ,$$
onde, sob a hipótese nula e condicionalmente a $\boldsymbol{X}$, $\frac{N}{D}$ tem a distribuição de uma razão de qui-quadrados independentes, com $r$ e $n-k$ graus de liberdade respectivamente. Portanto, pela definição da distribuição $F$, sob a nula e condicionalmente a $\boldsymbol{X}$, $W_{R,c}$ tem distribuição $F$ com $(r,n-k)$ graus de liberdade.


	


	
\end{document}